% =============================================================================
% 019 / 068
% Status: raw
% =============================================================================
% Notes:
%
% Source question:
% 您老有計劃做原生吳語的比較系統的發音、文法教學影片嗎？
% =============================================================================

\chapter[您老有計劃做原生吳語的比較系統的發音、文法教學影片嗎？]{您老有計劃做原生吳語的比較系統的發音、文法教學影片嗎？}
\qaanswer{
文法一直有出啊，另外吳語沒有所謂標準音，所以吳語拼音這種一個音位對應多個音值的拼音方案是最佳選。人為定標準音那是大一統帝國思維，都2019年了，觀念是不是得轉變一下？經驗主義的思路是使用某種語言的有教養的人群會自然形成一種語音腔調方式，這種腔調就會成為該語言的標準音。吳語的高階化路還很漫長。
}

